% ============================================================
%  Conclusion Générale — Mémoire de Master BIBDA
%  CHOUFA Zouhair
% ============================================================
\chapter*{Conclusion Générale}
\addcontentsline{toc}{chapter}{Conclusion Générale}
\label{chap:conclusion_generale}

\section*{Synthèse du travail réalisé}
\addcontentsline{toc}{section}{Synthèse du travail réalisé}

Ce mémoire s'est donné pour objectif de répondre à une question posée
dès l'introduction générale~: comment concevoir un pipeline
automatisé, modulaire et reproductible, capable de transformer des
données tabulaires hétérogènes et multi-sources en un Knowledge Graph
unifié, tout en garantissant la qualité et la validité formelle du
graphe produit à chaque étape du processus~? Le cheminement suivi
pour y répondre s'est appuyé sur quatre chapitres complémentaires.

Le \textbf{Chapitre~1} a dressé un état de l'art structuré, mettant
en évidence que le défi de la Variété au sein du Big Data ne pouvait
être résolu de façon automatisée et scalable par les approches
relationnelles classiques, et que si les fondements des Knowledge
Graphs (RDF, OWL, SPARQL, SHACL) offrent un cadre formel adapté, les
pipelines existants demeurent fragmentés et insuffisamment évalués.
Le \textbf{Chapitre~2} a traduit ce constat en une architecture
concrète, articulée en quatre modules~: ingestion et profilage,
\textit{schema matching} hybride, triplification déclarative et
résolution d'entités, complétée par un protocole d'évaluation rigoureux
fondé sur un gold standard annoté. Le \textbf{Chapitre~3} a validé
cette architecture par une implémentation effective dans un
environnement Cloud reproductible, produisant des résultats
opérationnels concrets pour chacun des modules. Enfin, le
\textbf{Chapitre~4} a dépassé la simple description de ces résultats
pour proposer une évaluation critique et comparative, calibrant
explicitement les hyperparamètres du pipeline et le positionnant par
rapport à la littérature.

\section*{Principaux résultats obtenus}
\addcontentsline{toc}{section}{Principaux résultats obtenus}

Trois résultats principaux ressortent de ce travail. Premièrement,
l'architecture hybride de \textit{schema matching}, combinant
embeddings Sentence-BERT et agent LLM (LLaMA~3.1 via LangChain),
atteint un F1-score de 0,870 après calibration explicite du seuil de
décision ($\tau^{*} = 0{,}82$), démontrant empiriquement qu'une telle
combinaison égalise ou dépasse chacune de ses composantes utilisée
isolément, tout en maîtrisant le coût des appels API — un résultat
qui n'est pas un gain de précision arbitraire mais le fruit d'une
procédure d'optimisation reproductible. Deuxièmement, le module de
triplification sémantique, fondé sur le moteur Morph-KGC et des
règles YARRRML auto-générées, a permis de produire 82\,420 triplets
RDF en seulement 7~secondes, illustrant la scalabilité de l'approche
déclarative retenue. Troisièmement, le module de résolution
d'entités, évalué sur 1\,600 paires réelles issues des deux sources
simulées, atteint un F1-score de 0,990 après recalibrage des règles
de blocage ($\lambda^{*} = 0{,}90$), confirmant la robustesse du
modèle probabiliste de Fellegi-Sunter pour la génération des ponts
sémantiques \texttt{owl:sameAs} constituant le graphe unifié final.

Au-delà de ces métriques, la discussion critique menée au
Chapitre~4 autour des quatre menaces classiques à la validité
expérimentale — de construit, interne, externe et de conclusion — a
permis de circonscrire précisément les conditions de généralisation
de ces résultats, renforçant ainsi la rigueur scientifique de
l'évaluation proposée plutôt que de se limiter à un rapport de
performances brutes.

\section*{Limites du travail}
\addcontentsline{toc}{section}{Limites du travail}

Ce travail présente plusieurs limites qu'il convient de reconnaître
explicitement. Le pipeline a été évalué sur un jeu de données
entièrement \textbf{synthétique}, généré via la bibliothèque
\texttt{Faker} à partir d'un dataset transactionnel public~; si ce
choix garantit la maîtrise et la reproductibilité de la vérité
terrain, il limite la portée externe des résultats, les données
réelles présentant typiquement des irrégularités et incohérences
plus complexes. Par ailleurs, le modèle Splink de résolution
d'entités s'est révélé fortement \textbf{sensible aux règles de
blocage} retenues, un blocage mal calibré pouvant rendre le modèle
inopérant indépendamment du seuil de décision~$\lambda$. L'architecture
proposée suppose également l'existence préalable d'une
\textbf{ontologie de domaine} bien formée, ce qui n'est pas
systématiquement le cas dans tous les contextes industriels. Enfin,
à très grande échelle, le \textbf{coût et la latence des appels LLM}
distants pourraient constituer un goulot d'étranglement pour des
pipelines traitant simultanément un grand nombre de sources.

\section*{Perspectives de recherche}
\addcontentsline{toc}{section}{Perspectives de recherche}

Plusieurs prolongements de ce travail apparaissent naturellement à
l'issue de cette analyse. Une \textbf{validation sur des données
réelles anonymisées}, issues d'un contexte industriel effectif,
permettrait de confirmer la généralisabilité des performances
obtenues sur données synthétiques. La définition de \textbf{clés de
blocage adaptatives}, apprises automatiquement plutôt que fixées
manuellement, constitue une direction de recherche ouverte pour
renforcer la robustesse du module de résolution d'entités. L'ajout
d'une étape de \textbf{construction ontologique automatique} en
amont du pipeline permettrait de s'affranchir de l'hypothèse d'une
ontologie de domaine préexistante. Enfin, la \textbf{substitution des
appels LLM distants par des modèles open source déployés localement}
(par exemple Mistral ou LLaMA via Ollama), directement supportée par
l'abstraction LangChain sans modification de l'architecture du
Module~2, permettrait de réduire la latence cumulée et la dépendance
à des API tierces pour un déploiement à très grande échelle.

En définitive, ce mémoire démontre qu'il est possible de concevoir
un pipeline hybride, à la fois symbolique et neuronal, pour la
construction automatisée de Knowledge Graphs à partir de données
tabulaires hétérogènes, tout en soumettant chacune de ses composantes
à une évaluation quantitative rigoureuse. Il constitue, à ce titre,
une base solide et reproductible pour de futurs travaux visant
l'industrialisation de tels pipelines dans des contextes
organisationnels réels.